% Betriebszustände und Übergänge der Umgebung (abstrakt, gilt für Simulation und Realbetrieb).
% Eigene Darstellung (Art E, Kap. 6.3.1). Keine Quelle in der Caption. Fasst vier gekoppelte Zustandsmaschinen
% des Codes (Onboard-Regler am Roboter, Ersatzregler der Simulation, Armregler, Bedienmodus) zu einem Diagramm zusammen.
% Basis, je Zustand/Übergang (G1@25d5c1a, Pfade relativ zu g1pilot/):
%   Gedämpft -> Bereit ("Start"): docs/50-A-Echter-Roboter.md F Schritt 2 (Start quittiert und stellt aus Damp auf
%     Preparation); docs/30-T-Stehen-und-Gehen.md Zustandsmaschine (/g1pilot/start(True) -> HOLD, Sim) und
%     Tabelle loco_client (/g1pilot/start).
%   Bereit -> Balancierend (Bedienmodus wählen): docs/30-T-Stehen-und-Gehen.md (start_balancing -> STAND, start_walking
%     -> _enter_walk; real: Tabelle loco_client, Bodenkontakt vor "Walk", knee_load_min_nm); docs/10-A-Demo-GUI.md
%     Abschnitt 1 (genau eine Kachel aktiv); docs/50-A-Echter-Roboter.md E (GREIFEN/LAUFEN nur aus Preparation).
%   GREIFEN <-> GEHEN: docs/30-T-Stehen-und-Gehen.md (START BALANCING im WALK: Ausbremsen, dann STAND; START WALKING).
%   Totmann losgelassen: docs/51-A-Notaus.md Stufen ("Gehen stoppen"); docs/30-T loco_client (Deadman/Timeout).
%   Not-Halt von jedem Zustand, Wirkung: docs/51-T-Notaus.md "Wer was tut" (arm_controller, loco_client, loco_sim,
%     Inspire-Bridge); Beine bleiben stehen (real), Simulation: DAMP (docs/51-A-Notaus.md "In der Simulation").
%   Quittieren: docs/51-A-Notaus.md "Quittieren" (Gamepad: geht nicht, absichtlich; danach Arme schlaff bis GREIFEN).
%   Voll-Notaus -> Gedämpft: docs/51-T-Notaus.md (Damp zuerst, nach Voll-Notaus bleibt Damp), docs/51-A-Notaus.md.
%   Sturz erkannt -> Gedämpft (nur Simulation): docs/30-T-Stehen-und-Gehen.md Zustandsmaschine (_fallen -> DAMP).
% Nicht gezeichnet (nicht für beide Welten belegt oder Detail der Tabelle): Damp per Startleiste aus dem Stehen
%   (nur real), Abfangen und Übergabe WALK->STAND (nur Simulation), zweite Bestätigung vor "Walk" (nur real),
%   Sperre bei Hüfte nicht in der Mitte, Fernbedienung des Herstellers.
% Stil: ba-grafik.tex. Einbinden: % Betriebszustände und Übergänge der Umgebung (abstrakt, gilt für Simulation und Realbetrieb).
% Eigene Darstellung (Art E, Kap. 6.3.1). Keine Quelle in der Caption. Fasst vier gekoppelte Zustandsmaschinen
% des Codes (Onboard-Regler am Roboter, Ersatzregler der Simulation, Armregler, Bedienmodus) zu einem Diagramm zusammen.
% Basis, je Zustand/Übergang (G1@25d5c1a, Pfade relativ zu g1pilot/):
%   Gedämpft -> Bereit ("Start"): docs/50-A-Echter-Roboter.md F Schritt 2 (Start quittiert und stellt aus Damp auf
%     Preparation); docs/30-T-Stehen-und-Gehen.md Zustandsmaschine (/g1pilot/start(True) -> HOLD, Sim) und
%     Tabelle loco_client (/g1pilot/start).
%   Bereit -> Balancierend (Bedienmodus wählen): docs/30-T-Stehen-und-Gehen.md (start_balancing -> STAND, start_walking
%     -> _enter_walk; real: Tabelle loco_client, Bodenkontakt vor "Walk", knee_load_min_nm); docs/10-A-Demo-GUI.md
%     Abschnitt 1 (genau eine Kachel aktiv); docs/50-A-Echter-Roboter.md E (GREIFEN/LAUFEN nur aus Preparation).
%   GREIFEN <-> GEHEN: docs/30-T-Stehen-und-Gehen.md (START BALANCING im WALK: Ausbremsen, dann STAND; START WALKING).
%   Totmann losgelassen: docs/51-A-Notaus.md Stufen ("Gehen stoppen"); docs/30-T loco_client (Deadman/Timeout).
%   Not-Halt von jedem Zustand, Wirkung: docs/51-T-Notaus.md "Wer was tut" (arm_controller, loco_client, loco_sim,
%     Inspire-Bridge); Beine bleiben stehen (real), Simulation: DAMP (docs/51-A-Notaus.md "In der Simulation").
%   Quittieren: docs/51-A-Notaus.md "Quittieren" (Gamepad: geht nicht, absichtlich; danach Arme schlaff bis GREIFEN).
%   Voll-Notaus -> Gedämpft: docs/51-T-Notaus.md (Damp zuerst, nach Voll-Notaus bleibt Damp), docs/51-A-Notaus.md.
%   Sturz erkannt -> Gedämpft (nur Simulation): docs/30-T-Stehen-und-Gehen.md Zustandsmaschine (_fallen -> DAMP).
% Nicht gezeichnet (nicht für beide Welten belegt oder Detail der Tabelle): Damp per Startleiste aus dem Stehen
%   (nur real), Abfangen und Übergabe WALK->STAND (nur Simulation), zweite Bestätigung vor "Walk" (nur real),
%   Sperre bei Hüfte nicht in der Mitte, Fernbedienung des Herstellers.
% Stil: ba-grafik.tex. Einbinden: % Betriebszustände und Übergänge der Umgebung (abstrakt, gilt für Simulation und Realbetrieb).
% Eigene Darstellung (Art E, Kap. 6.3.1). Keine Quelle in der Caption. Fasst vier gekoppelte Zustandsmaschinen
% des Codes (Onboard-Regler am Roboter, Ersatzregler der Simulation, Armregler, Bedienmodus) zu einem Diagramm zusammen.
% Basis, je Zustand/Übergang (G1@25d5c1a, Pfade relativ zu g1pilot/):
%   Gedämpft -> Bereit ("Start"): docs/50-A-Echter-Roboter.md F Schritt 2 (Start quittiert und stellt aus Damp auf
%     Preparation); docs/30-T-Stehen-und-Gehen.md Zustandsmaschine (/g1pilot/start(True) -> HOLD, Sim) und
%     Tabelle loco_client (/g1pilot/start).
%   Bereit -> Balancierend (Bedienmodus wählen): docs/30-T-Stehen-und-Gehen.md (start_balancing -> STAND, start_walking
%     -> _enter_walk; real: Tabelle loco_client, Bodenkontakt vor "Walk", knee_load_min_nm); docs/10-A-Demo-GUI.md
%     Abschnitt 1 (genau eine Kachel aktiv); docs/50-A-Echter-Roboter.md E (GREIFEN/LAUFEN nur aus Preparation).
%   GREIFEN <-> GEHEN: docs/30-T-Stehen-und-Gehen.md (START BALANCING im WALK: Ausbremsen, dann STAND; START WALKING).
%   Totmann losgelassen: docs/51-A-Notaus.md Stufen ("Gehen stoppen"); docs/30-T loco_client (Deadman/Timeout).
%   Not-Halt von jedem Zustand, Wirkung: docs/51-T-Notaus.md "Wer was tut" (arm_controller, loco_client, loco_sim,
%     Inspire-Bridge); Beine bleiben stehen (real), Simulation: DAMP (docs/51-A-Notaus.md "In der Simulation").
%   Quittieren: docs/51-A-Notaus.md "Quittieren" (Gamepad: geht nicht, absichtlich; danach Arme schlaff bis GREIFEN).
%   Voll-Notaus -> Gedämpft: docs/51-T-Notaus.md (Damp zuerst, nach Voll-Notaus bleibt Damp), docs/51-A-Notaus.md.
%   Sturz erkannt -> Gedämpft (nur Simulation): docs/30-T-Stehen-und-Gehen.md Zustandsmaschine (_fallen -> DAMP).
% Nicht gezeichnet (nicht für beide Welten belegt oder Detail der Tabelle): Damp per Startleiste aus dem Stehen
%   (nur real), Abfangen und Übergabe WALK->STAND (nur Simulation), zweite Bestätigung vor "Walk" (nur real),
%   Sperre bei Hüfte nicht in der Mitte, Fernbedienung des Herstellers.
% Stil: ba-grafik.tex. Einbinden: % Betriebszustände und Übergänge der Umgebung (abstrakt, gilt für Simulation und Realbetrieb).
% Eigene Darstellung (Art E, Kap. 6.3.1). Keine Quelle in der Caption. Fasst vier gekoppelte Zustandsmaschinen
% des Codes (Onboard-Regler am Roboter, Ersatzregler der Simulation, Armregler, Bedienmodus) zu einem Diagramm zusammen.
% Basis, je Zustand/Übergang (G1@25d5c1a, Pfade relativ zu g1pilot/):
%   Gedämpft -> Bereit ("Start"): docs/50-A-Echter-Roboter.md F Schritt 2 (Start quittiert und stellt aus Damp auf
%     Preparation); docs/30-T-Stehen-und-Gehen.md Zustandsmaschine (/g1pilot/start(True) -> HOLD, Sim) und
%     Tabelle loco_client (/g1pilot/start).
%   Bereit -> Balancierend (Bedienmodus wählen): docs/30-T-Stehen-und-Gehen.md (start_balancing -> STAND, start_walking
%     -> _enter_walk; real: Tabelle loco_client, Bodenkontakt vor "Walk", knee_load_min_nm); docs/10-A-Demo-GUI.md
%     Abschnitt 1 (genau eine Kachel aktiv); docs/50-A-Echter-Roboter.md E (GREIFEN/LAUFEN nur aus Preparation).
%   GREIFEN <-> GEHEN: docs/30-T-Stehen-und-Gehen.md (START BALANCING im WALK: Ausbremsen, dann STAND; START WALKING).
%   Totmann losgelassen: docs/51-A-Notaus.md Stufen ("Gehen stoppen"); docs/30-T loco_client (Deadman/Timeout).
%   Not-Halt von jedem Zustand, Wirkung: docs/51-T-Notaus.md "Wer was tut" (arm_controller, loco_client, loco_sim,
%     Inspire-Bridge); Beine bleiben stehen (real), Simulation: DAMP (docs/51-A-Notaus.md "In der Simulation").
%   Quittieren: docs/51-A-Notaus.md "Quittieren" (Gamepad: geht nicht, absichtlich; danach Arme schlaff bis GREIFEN).
%   Voll-Notaus -> Gedämpft: docs/51-T-Notaus.md (Damp zuerst, nach Voll-Notaus bleibt Damp), docs/51-A-Notaus.md.
%   Sturz erkannt -> Gedämpft (nur Simulation): docs/30-T-Stehen-und-Gehen.md Zustandsmaschine (_fallen -> DAMP).
% Nicht gezeichnet (nicht für beide Welten belegt oder Detail der Tabelle): Damp per Startleiste aus dem Stehen
%   (nur real), Abfangen und Übergabe WALK->STAND (nur Simulation), zweite Bestätigung vor "Walk" (nur real),
%   Sperre bei Hüfte nicht in der Mitte, Fernbedienung des Herstellers.
% Stil: ba-grafik.tex. Einbinden: \input{Bilder/tikz/betriebszustaende} in einer figure-Umgebung. Nicht skalieren.
\begin{tikzpicture}[ba]
  \tikzset{
    zu/.style={ba block, align=center, minimum height=12mm},
    halt/.style={ba block, draw=baAkzent, line width=1.2pt, align=center, minimum height=12mm},
    ueb/.style={font=\footnotesize, text=black},
    klein/.style={font=\footnotesize, text=baGrau},
  }
  % Rahmen: Not-Halt und Voll-Notaus gelten von jedem Zustand
  \draw[ba gruppe] (0.3,1.85) rectangle (15.7,8.95);
  \node[klein, anchor=south west] at (0.45,1.95) {von jedem Zustand};
  % Zustände
  \node[zu, text width=30mm] (G) at (2.3,7.9) {Gedämpft\\{\footnotesize alle Gelenke weich}};
  \node[zu, text width=44mm] (B) at (8.2,7.9) {Bereit\\{\footnotesize Gelenke gehalten, noch kein Balancieren}};
  \draw[ba gruppe, draw=baGrau] (4.0,2.3) rectangle (15.3,5.7);
  \node[font=\small, anchor=north west] at (4.15,5.65) {Balancierend (genau ein Bedienmodus aktiv)};
  \node[zu, text width=34mm, ba block hervor] (GR) at (6.1,3.95) {GREIFEN\\{\footnotesize Stehen, Arme frei}};
  \node[zu, text width=40mm, ba block hervor] (GE) at (13.0,3.95) {GEHEN\\{\footnotesize Geschwindigkeitsvorgabe}};
  % Übergänge
  \draw[ba pfeil] (G.east) -- node[ueb, above] {Start} (B.west);
  \draw[ba pfeil] (B.south) -- (8.2,5.7);
  \node[ueb, anchor=west] at (8.4,6.45) {Modus wählen (real: nur mit Bodenkontakt)};
  \draw[ba pfeil] ([yshift=2mm]GR.east) -- node[ueb, above] {GEHEN wählen} ([yshift=2mm]GE.west);
  \draw[ba pfeil] ([yshift=-2mm]GE.west) -- node[ueb, below] {GREIFEN wählen} ([yshift=-2mm]GR.east);
  \node[klein, anchor=south, align=center] at (13.0,2.38) {Totmann losgelassen:\\Gehen stoppt};
  % Sturz (nur Simulation)
  \draw[ba pfeil] (4.0,3.95) -| ([xshift=8mm]G.south);
  \node[ueb, anchor=west, align=left] at (3.2,6.45) {Sturz erkannt\\(nur Simulation)};
  % Not-Halt (Sammelübergang vom Rahmen)
  \node[halt, text width=50mm] (NH) at (6.5,0.1) {Not-Halt (verriegelt)\\{\footnotesize Arme schlaff, Hände halten, Taille gehalten, Beine bleiben stehen}};
  \draw[ba pfeil, baAkzent] (6.5,1.85) -- (NH.north);
  \node[ueb, anchor=west, align=left] at (6.65,1.4) {auf jedem\\Bedienweg};
  % Quittieren
  \draw[ba pfeil] (NH.east) -| (14.5,2.3);
  \node[ueb, anchor=south west, text width=47mm] at (9.6,0.35) {Quittieren (nicht am \mbox{Gamepad}); Arme bleiben schlaff, bis GREIFEN gewählt wird};
  % Voll-Notaus
  \node[halt, text width=24mm] (VN) at (2.1,0.1) {Voll-Notaus\\{\footnotesize alles gedämpft}};
  \draw[ba pfeil, baAkzent] (2.1,1.85) -- (VN.north);
  \node[ueb, anchor=west, text width=30mm] at (2.25,1.4) {Gamepad, Tasten gehalten};
  \draw[ba pfeil] (VN.west) -- ++(-0.2,0) |- (G.west);
\end{tikzpicture}
 in einer figure-Umgebung. Nicht skalieren.
\begin{tikzpicture}[ba]
  \tikzset{
    zu/.style={ba block, align=center, minimum height=12mm},
    halt/.style={ba block, draw=baAkzent, line width=1.2pt, align=center, minimum height=12mm},
    ueb/.style={font=\footnotesize, text=black},
    klein/.style={font=\footnotesize, text=baGrau},
  }
  % Rahmen: Not-Halt und Voll-Notaus gelten von jedem Zustand
  \draw[ba gruppe] (0.3,1.85) rectangle (15.7,8.95);
  \node[klein, anchor=south west] at (0.45,1.95) {von jedem Zustand};
  % Zustände
  \node[zu, text width=30mm] (G) at (2.3,7.9) {Gedämpft\\{\footnotesize alle Gelenke weich}};
  \node[zu, text width=44mm] (B) at (8.2,7.9) {Bereit\\{\footnotesize Gelenke gehalten, noch kein Balancieren}};
  \draw[ba gruppe, draw=baGrau] (4.0,2.3) rectangle (15.3,5.7);
  \node[font=\small, anchor=north west] at (4.15,5.65) {Balancierend (genau ein Bedienmodus aktiv)};
  \node[zu, text width=34mm, ba block hervor] (GR) at (6.1,3.95) {GREIFEN\\{\footnotesize Stehen, Arme frei}};
  \node[zu, text width=40mm, ba block hervor] (GE) at (13.0,3.95) {GEHEN\\{\footnotesize Geschwindigkeitsvorgabe}};
  % Übergänge
  \draw[ba pfeil] (G.east) -- node[ueb, above] {Start} (B.west);
  \draw[ba pfeil] (B.south) -- (8.2,5.7);
  \node[ueb, anchor=west] at (8.4,6.45) {Modus wählen (real: nur mit Bodenkontakt)};
  \draw[ba pfeil] ([yshift=2mm]GR.east) -- node[ueb, above] {GEHEN wählen} ([yshift=2mm]GE.west);
  \draw[ba pfeil] ([yshift=-2mm]GE.west) -- node[ueb, below] {GREIFEN wählen} ([yshift=-2mm]GR.east);
  \node[klein, anchor=south, align=center] at (13.0,2.38) {Totmann losgelassen:\\Gehen stoppt};
  % Sturz (nur Simulation)
  \draw[ba pfeil] (4.0,3.95) -| ([xshift=8mm]G.south);
  \node[ueb, anchor=west, align=left] at (3.2,6.45) {Sturz erkannt\\(nur Simulation)};
  % Not-Halt (Sammelübergang vom Rahmen)
  \node[halt, text width=50mm] (NH) at (6.5,0.1) {Not-Halt (verriegelt)\\{\footnotesize Arme schlaff, Hände halten, Taille gehalten, Beine bleiben stehen}};
  \draw[ba pfeil, baAkzent] (6.5,1.85) -- (NH.north);
  \node[ueb, anchor=west, align=left] at (6.65,1.4) {auf jedem\\Bedienweg};
  % Quittieren
  \draw[ba pfeil] (NH.east) -| (14.5,2.3);
  \node[ueb, anchor=south west, text width=47mm] at (9.6,0.35) {Quittieren (nicht am \mbox{Gamepad}); Arme bleiben schlaff, bis GREIFEN gewählt wird};
  % Voll-Notaus
  \node[halt, text width=24mm] (VN) at (2.1,0.1) {Voll-Notaus\\{\footnotesize alles gedämpft}};
  \draw[ba pfeil, baAkzent] (2.1,1.85) -- (VN.north);
  \node[ueb, anchor=west, text width=30mm] at (2.25,1.4) {Gamepad, Tasten gehalten};
  \draw[ba pfeil] (VN.west) -- ++(-0.2,0) |- (G.west);
\end{tikzpicture}
 in einer figure-Umgebung. Nicht skalieren.
\begin{tikzpicture}[ba]
  \tikzset{
    zu/.style={ba block, align=center, minimum height=12mm},
    halt/.style={ba block, draw=baAkzent, line width=1.2pt, align=center, minimum height=12mm},
    ueb/.style={font=\footnotesize, text=black},
    klein/.style={font=\footnotesize, text=baGrau},
  }
  % Rahmen: Not-Halt und Voll-Notaus gelten von jedem Zustand
  \draw[ba gruppe] (0.3,1.85) rectangle (15.7,8.95);
  \node[klein, anchor=south west] at (0.45,1.95) {von jedem Zustand};
  % Zustände
  \node[zu, text width=30mm] (G) at (2.3,7.9) {Gedämpft\\{\footnotesize alle Gelenke weich}};
  \node[zu, text width=44mm] (B) at (8.2,7.9) {Bereit\\{\footnotesize Gelenke gehalten, noch kein Balancieren}};
  \draw[ba gruppe, draw=baGrau] (4.0,2.3) rectangle (15.3,5.7);
  \node[font=\small, anchor=north west] at (4.15,5.65) {Balancierend (genau ein Bedienmodus aktiv)};
  \node[zu, text width=34mm, ba block hervor] (GR) at (6.1,3.95) {GREIFEN\\{\footnotesize Stehen, Arme frei}};
  \node[zu, text width=40mm, ba block hervor] (GE) at (13.0,3.95) {GEHEN\\{\footnotesize Geschwindigkeitsvorgabe}};
  % Übergänge
  \draw[ba pfeil] (G.east) -- node[ueb, above] {Start} (B.west);
  \draw[ba pfeil] (B.south) -- (8.2,5.7);
  \node[ueb, anchor=west] at (8.4,6.45) {Modus wählen (real: nur mit Bodenkontakt)};
  \draw[ba pfeil] ([yshift=2mm]GR.east) -- node[ueb, above] {GEHEN wählen} ([yshift=2mm]GE.west);
  \draw[ba pfeil] ([yshift=-2mm]GE.west) -- node[ueb, below] {GREIFEN wählen} ([yshift=-2mm]GR.east);
  \node[klein, anchor=south, align=center] at (13.0,2.38) {Totmann losgelassen:\\Gehen stoppt};
  % Sturz (nur Simulation)
  \draw[ba pfeil] (4.0,3.95) -| ([xshift=8mm]G.south);
  \node[ueb, anchor=west, align=left] at (3.2,6.45) {Sturz erkannt\\(nur Simulation)};
  % Not-Halt (Sammelübergang vom Rahmen)
  \node[halt, text width=50mm] (NH) at (6.5,0.1) {Not-Halt (verriegelt)\\{\footnotesize Arme schlaff, Hände halten, Taille gehalten, Beine bleiben stehen}};
  \draw[ba pfeil, baAkzent] (6.5,1.85) -- (NH.north);
  \node[ueb, anchor=west, align=left] at (6.65,1.4) {auf jedem\\Bedienweg};
  % Quittieren
  \draw[ba pfeil] (NH.east) -| (14.5,2.3);
  \node[ueb, anchor=south west, text width=47mm] at (9.6,0.35) {Quittieren (nicht am \mbox{Gamepad}); Arme bleiben schlaff, bis GREIFEN gewählt wird};
  % Voll-Notaus
  \node[halt, text width=24mm] (VN) at (2.1,0.1) {Voll-Notaus\\{\footnotesize alles gedämpft}};
  \draw[ba pfeil, baAkzent] (2.1,1.85) -- (VN.north);
  \node[ueb, anchor=west, text width=30mm] at (2.25,1.4) {Gamepad, Tasten gehalten};
  \draw[ba pfeil] (VN.west) -- ++(-0.2,0) |- (G.west);
\end{tikzpicture}
 in einer figure-Umgebung. Nicht skalieren.
\begin{tikzpicture}[ba]
  \tikzset{
    zu/.style={ba block, align=center, minimum height=12mm},
    halt/.style={ba block, draw=baAkzent, line width=1.2pt, align=center, minimum height=12mm},
    ueb/.style={font=\footnotesize, text=black},
    klein/.style={font=\footnotesize, text=baGrau},
  }
  % Rahmen: Not-Halt und Voll-Notaus gelten von jedem Zustand
  \draw[ba gruppe] (0.3,1.85) rectangle (15.7,8.95);
  \node[klein, anchor=south west] at (0.45,1.95) {von jedem Zustand};
  % Zustände
  \node[zu, text width=30mm] (G) at (2.3,7.9) {Gedämpft\\{\footnotesize alle Gelenke weich}};
  \node[zu, text width=44mm] (B) at (8.2,7.9) {Bereit\\{\footnotesize Gelenke gehalten, noch kein Balancieren}};
  \draw[ba gruppe, draw=baGrau] (4.0,2.3) rectangle (15.3,5.7);
  \node[font=\small, anchor=north west] at (4.15,5.65) {Balancierend (genau ein Bedienmodus aktiv)};
  \node[zu, text width=34mm, ba block hervor] (GR) at (6.1,3.95) {GREIFEN\\{\footnotesize Stehen, Arme frei}};
  \node[zu, text width=40mm, ba block hervor] (GE) at (13.0,3.95) {GEHEN\\{\footnotesize Geschwindigkeitsvorgabe}};
  % Übergänge
  \draw[ba pfeil] (G.east) -- node[ueb, above] {Start} (B.west);
  \draw[ba pfeil] (B.south) -- (8.2,5.7);
  \node[ueb, anchor=west] at (8.4,6.45) {Modus wählen (real: nur mit Bodenkontakt)};
  \draw[ba pfeil] ([yshift=2mm]GR.east) -- node[ueb, above] {GEHEN wählen} ([yshift=2mm]GE.west);
  \draw[ba pfeil] ([yshift=-2mm]GE.west) -- node[ueb, below] {GREIFEN wählen} ([yshift=-2mm]GR.east);
  \node[klein, anchor=south, align=center] at (13.0,2.38) {Totmann losgelassen:\\Gehen stoppt};
  % Sturz (nur Simulation)
  \draw[ba pfeil] (4.0,3.95) -| ([xshift=8mm]G.south);
  \node[ueb, anchor=west, align=left] at (3.2,6.45) {Sturz erkannt\\(nur Simulation)};
  % Not-Halt (Sammelübergang vom Rahmen)
  \node[halt, text width=50mm] (NH) at (6.5,0.1) {Not-Halt (verriegelt)\\{\footnotesize Arme schlaff, Hände halten, Taille gehalten, Beine bleiben stehen}};
  \draw[ba pfeil, baAkzent] (6.5,1.85) -- (NH.north);
  \node[ueb, anchor=west, align=left] at (6.65,1.4) {auf jedem\\Bedienweg};
  % Quittieren
  \draw[ba pfeil] (NH.east) -| (14.5,2.3);
  \node[ueb, anchor=south west, text width=47mm] at (9.6,0.35) {Quittieren (nicht am \mbox{Gamepad}); Arme bleiben schlaff, bis GREIFEN gewählt wird};
  % Voll-Notaus
  \node[halt, text width=24mm] (VN) at (2.1,0.1) {Voll-Notaus\\{\footnotesize alles gedämpft}};
  \draw[ba pfeil, baAkzent] (2.1,1.85) -- (VN.north);
  \node[ueb, anchor=west, text width=30mm] at (2.25,1.4) {Gamepad, Tasten gehalten};
  \draw[ba pfeil] (VN.west) -- ++(-0.2,0) |- (G.west);
\end{tikzpicture}
